\documentclass[a4paper,12pt,twoside]{memoir}

% Castellano
\usepackage[spanish,es-tabla]{babel}
\selectlanguage{spanish}
\usepackage[utf8]{inputenc}
\usepackage[T1]{fontenc}
\usepackage{lmodern} % Scalable font
\usepackage{microtype}
\usepackage{placeins}

\RequirePackage{booktabs}
\RequirePackage[table]{xcolor}
\RequirePackage{xtab}
\RequirePackage{multirow}

% Links
\PassOptionsToPackage{hyphens}{url}\usepackage[colorlinks]{hyperref}
\hypersetup{
	allcolors = {red}
}

% Ecuaciones
\usepackage{amsmath}

% Rutas de fichero / paquete
\newcommand{\ruta}[1]{{\sffamily #1}}

% Párrafos
\nonzeroparskip

% Huérfanas y viudas
\widowpenalty100000
\clubpenalty100000

% Imágenes

% Comando para insertar una imagen en un lugar concreto.
% Los parámetros son:
% 1 --> Ruta absoluta/relativa de la figura
% 2 --> Texto a pie de figura
% 3 --> Tamaño en tanto por uno relativo al ancho de página
\usepackage{graphicx}
\newcommand{\imagen}[3]{
	\begin{figure}[!h]
		\centering
		\includegraphics[width=#3\textwidth]{#1}
		\caption{#2}\label{fig:#1}
	\end{figure}
	\FloatBarrier
}

% Comando para insertar una imagen sin posición.
% Los parámetros son:
% 1 --> Ruta absoluta/relativa de la figura
% 2 --> Texto a pie de figura
% 3 --> Tamaño en tanto por uno relativo al ancho de página
\newcommand{\imagenflotante}[3]{
	\begin{figure}
		\centering
		\includegraphics[width=#3\textwidth]{#1}
		\caption{#2}\label{fig:#1}
	\end{figure}
}

% El comando \figura nos permite insertar figuras comodamente, y utilizando
% siempre el mismo formato. Los parametros son:
% 1 --> Porcentaje del ancho de página que ocupará la figura (de 0 a 1)
% 2 --> Fichero de la imagen
% 3 --> Texto a pie de imagen
% 4 --> Etiqueta (label) para referencias
% 5 --> Opciones que queramos pasarle al \includegraphics
% 6 --> Opciones de posicionamiento a pasarle a \begin{figure}
\newcommand{\figuraConPosicion}[6]{%
  \setlength{\anchoFloat}{#1\textwidth}%
  \addtolength{\anchoFloat}{-4\fboxsep}%
  \setlength{\anchoFigura}{\anchoFloat}%
  \begin{figure}[#6]
    \begin{center}%
      \Ovalbox{%
        \begin{minipage}{\anchoFloat}%
          \begin{center}%
            \includegraphics[width=\anchoFigura,#5]{#2}%
            \caption{#3}%
            \label{#4}%
          \end{center}%
        \end{minipage}
      }%
    \end{center}%
  \end{figure}%
}

%
% Comando para incluir imágenes en formato apaisado (sin marco).
\newcommand{\figuraApaisadaSinMarco}[5]{%
  \begin{figure}%
    \begin{center}%
    \includegraphics[angle=90,height=#1\textheight,#5]{#2}%
    \caption{#3}%
    \label{#4}%
    \end{center}%
  \end{figure}%
}
% Para las tablas
\newcommand{\otoprule}{\midrule [\heavyrulewidth]}
%
% Nuevo comando para tablas pequeñas (menos de una página).
\newcommand{\tablaSmall}[5]{%
 \begin{table}
  \begin{center}
   \rowcolors {2}{gray!35}{}
   \begin{tabular}{#2}
    \toprule
    #4
    \otoprule
    #5
    \bottomrule
   \end{tabular}
   \caption{#1}
   \label{tabla:#3}
  \end{center}
 \end{table}
}

%
% Nuevo comando para tablas pequeñas (menos de una página).
\newcommand{\tablaSmallSinColores}[5]{%
 \begin{table}[H]
  \begin{center}
   \begin{tabular}{#2}
    \toprule
    #4
    \otoprule
    #5
    \bottomrule
   \end{tabular}
   \caption{#1}
   \label{tabla:#3}
  \end{center}
 \end{table}
}

\newcommand{\tablaApaisadaSmall}[5]{%
\begin{landscape}
  \begin{table}
   \begin{center}
    \rowcolors {2}{gray!35}{}
    \begin{tabular}{#2}
     \toprule
     #4
     \otoprule
     #5
     \bottomrule
    \end{tabular}
    \caption{#1}
    \label{tabla:#3}
   \end{center}
  \end{table}
\end{landscape}
}

%
% Nuevo comando para tablas grandes con cabecera y filas alternas coloreadas en gris.
\newcommand{\tabla}[6]{%
  \begin{center}
    \tablefirsthead{
      \toprule
      #5
      \otoprule
    }
    \tablehead{
      \multicolumn{#3}{l}{\small\sl continúa desde la página anterior}\\
      \toprule
      #5
      \otoprule
    }
    \tabletail{
      \hline
      \multicolumn{#3}{r}{\small\sl continúa en la página siguiente}\\
    }
    \tablelasttail{
      \hline
    }
    \bottomcaption{#1}
    \rowcolors {2}{gray!35}{}
    \begin{xtabular}{#2}
      #6
      \bottomrule
    \end{xtabular}
    \label{tabla:#4}
  \end{center}
}

%
% Nuevo comando para tablas grandes con cabecera.
\newcommand{\tablaSinColores}[6]{%
  \begin{center}
    \tablefirsthead{
      \toprule
      #5
      \otoprule
    }
    \tablehead{
      \multicolumn{#3}{l}{\small\sl continúa desde la página anterior}\\
      \toprule
      #5
      \otoprule
    }
    \tabletail{
      \hline
      \multicolumn{#3}{r}{\small\sl continúa en la página siguiente}\\
    }
    \tablelasttail{
      \hline
    }
    \bottomcaption{#1}
    \begin{xtabular}{#2}
      #6
      \bottomrule
    \end{xtabular}
    \label{tabla:#4}
  \end{center}
}

%
% Nuevo comando para tablas grandes sin cabecera.
\newcommand{\tablaSinCabecera}[5]{%
  \begin{center}
    \tablefirsthead{
      \toprule
    }
    \tablehead{
      \multicolumn{#3}{l}{\small\sl continúa desde la página anterior}\\
      \hline
    }
    \tabletail{
      \hline
      \multicolumn{#3}{r}{\small\sl continúa en la página siguiente}\\
    }
    \tablelasttail{
      \hline
    }
    \bottomcaption{#1}
  \begin{xtabular}{#2}
    #5
   \bottomrule
  \end{xtabular}
  \label{tabla:#4}
  \end{center}
}



\definecolor{cgoLight}{HTML}{EEEEEE}
\definecolor{cgoExtralight}{HTML}{FFFFFF}

%
% Nuevo comando para tablas grandes sin cabecera.
\newcommand{\tablaSinCabeceraConBandas}[5]{%
  \begin{center}
    \tablefirsthead{
      \toprule
    }
    \tablehead{
      \multicolumn{#3}{l}{\small\sl continúa desde la página anterior}\\
      \hline
    }
    \tabletail{
      \hline
      \multicolumn{#3}{r}{\small\sl continúa en la página siguiente}\\
    }
    \tablelasttail{
      \hline
    }
    \bottomcaption{#1}
    \rowcolors[]{1}{cgoExtralight}{cgoLight}

  \begin{xtabular}{#2}
    #5
   \bottomrule
  \end{xtabular}
  \label{tabla:#4}
  \end{center}
}



\graphicspath{ {./img/} }

% Capítulos
\chapterstyle{bianchi}
\newcommand{\capitulo}[2]{
	\setcounter{chapter}{#1}
	\setcounter{section}{0}
	\setcounter{figure}{0}
	\setcounter{table}{0}
	\chapter*{\thechapter.\enskip #2}
	\addcontentsline{toc}{chapter}{\thechapter.\enskip #2}
	\markboth{#2}{#2}
}

% Apéndices
\renewcommand{\appendixname}{Apéndice}
\renewcommand*\cftappendixname{\appendixname}

\newcommand{\apendice}[1]{
	%\renewcommand{\thechapter}{A}
	\chapter{#1}
}

\renewcommand*\cftappendixname{\appendixname\ }

% Formato de portada
\makeatletter
\usepackage{xcolor}
\newcommand{\tutor}[1]{\def\@tutor{#1}}
\newcommand{\course}[1]{\def\@course{#1}}
\definecolor{cpardoBox}{HTML}{E6E6FF}
\def\maketitle{
  \null
  \thispagestyle{empty}
  % Cabecera ----------------
\noindent\includegraphics[width=\textwidth]{cabecera}\vspace{1cm}%
  \vfill
  % Título proyecto y escudo informática ----------------
  \colorbox{cpardoBox}{%
    \begin{minipage}{.8\textwidth}
      \vspace{.5cm}\Large
      \begin{center}
      \textbf{TFG del Grado en Ingeniería Informática}\vspace{.6cm}\\
      \textbf{\LARGE\@title{}}
      \end{center}
      \vspace{.2cm}
    \end{minipage}

  }%
  \hfill\begin{minipage}{.20\textwidth}
    \includegraphics[width=\textwidth]{escudoInfor}
  \end{minipage}
  \vfill
  % Datos de alumno, curso y tutores ------------------
  \begin{center}%
  {%
    \noindent\LARGE
    Presentado por \@author{}\\ 
    en Universidad de Burgos --- \@date{}\\
    Tutor: \@tutor{}\\
  }%
  \end{center}%
  \null
  \cleardoublepage
  }
\makeatother

\newcommand{\nombre}{Rodrigo Pérez Ubierna} %%% cambio de comando

% Datos de portada
\title{{Kairos}\\[0.5cm]Aplicación de compra y venta de relojes por subasta.}
\author{\nombre}
\tutor{Dra. Sandra Rodríguez Arribas\\y Dr. Jose Antonio Barbero Aparicio}
\date{\today}

\begin{document}

\maketitle


\newpage\null\thispagestyle{empty}\newpage


%%%%%%%%%%%%%%%%%%%%%%%%%%%%%%%%%%%%%%%%%%%%%%%%%%%%%%%%%%%%%%%%%%%%%%%%%%%%%%%%%%%%%%%%
\thispagestyle{empty}


\noindent\includegraphics[width=\textwidth]{cabecera}\vspace{1cm}

\noindent Dña. Sandra Rodríguez Arribas y D. José Antonio Barbero Aparicio, profesores del Departamento de Ingeniería Informática, Área de Lenguajes y Sistemas Informáticos.

\noindent Exponen:

\noindent Que el alumno D. \nombre, con DNI 71708828A, ha realizado el Trabajo final de Grado en Ingeniería Informática titulado Kairos. 

\noindent Y que dicho trabajo ha sido realizado por el alumno bajo la dirección del que suscribe, en virtud de lo cual se autoriza su presentación y defensa.

\begin{center} %\large
En Burgos, {\large \today}
\end{center}

\vfill\vfill\vfill

% Author and supervisor
\begin{minipage}{0.45\textwidth}
\begin{flushleft} %\large
Vº. Bº. del Tutor:\\[2cm]
Dña. Sandra Rodríguez Arribas
\end{flushleft}
\end{minipage}
\hfill
\begin{minipage}{0.45\textwidth}
\begin{flushleft} %\large
Vº. Bº. del co-tutor:\\[2cm]
D. José Antonio Barbero Aparicio
\end{flushleft}
\end{minipage}
\hfill

\vfill

% para casos con solo un tutor comentar lo anterior
% y descomentar lo siguiente
%Vº. Bº. del Tutor:\\[2cm]
%D. nombre tutor


\newpage\null\thispagestyle{empty}\newpage




\frontmatter

% Abstract en castellano
\renewcommand*\abstractname{Resumen}
\begin{abstract}
El mundo de los relojes es un arte muy presente en la actualidad. Kairos nace por la necesidad de llevar este mundo a más gente. Kairos es una aplicación multiplataforma que permite la venta de relojes por subasta y/o venta directa. Siendo una aplicación fácil de utilizar y accesible a cualquier usuario, en ella se podrá registrar un usuario, subir los relojes y venderlos a través de la realización de una suabsta de tipo ascendente. Además, Kairos cuenta con la conexión a un modelo de predicción de precios de cualquier reloj, disponible a la hora de subir un reloj. La parte frontal de la aplicación se realiza con Flutter y estará conectada a una base de datos en la nube con el nombre Firestore Database.
\end{abstract}

\renewcommand*\abstractname{Descriptores}
\begin{abstract}
Relojes, aplicación multiplataforma, venta de relojes, subasta, venta directa, usuario, registro, predicción de precios, Flutter, Firestore Database, base de datos en la nube \ldots
\end{abstract}

\clearpage

% Abstract en inglés
\renewcommand*\abstractname{Abstract}
\begin{abstract}
The world of watches is a very present art nowadays. Kairos was born out of the need to bring this world to more people. Kairos is a multiplatform application that allows the sale of watches by auction and/or direct sale. Being an easy to use application and accessible to any user, it will be possible to register a user, upload the watches and sell them through an ascending auction. In addition, Kairos has a connection to a price prediction model of any watch, available when uploading a watch. The front end of the application is realised with Flutter and will be connected to a cloud database called Firestore Database.
\end{abstract}

\renewcommand*\abstractname{Keywords}
\begin{abstract}
Watches, cross-platform application, watch sale, auction, direct sale, user, registration, price prediction, Flutter, Firestore Database, cloud database.
\end{abstract}

\clearpage

% Indices
\tableofcontents

\clearpage

\listoffigures

\clearpage

\listoftables
\clearpage

\mainmatter
\capitulo{1}{Introducción}

	Vender un producto en internet está a la orden del día. Pero si les  Puede que algunos no, pero existe una gran mayoría de coleccionistas que invierten horas en conocer cuánto vale su producto. Esto es así porque lo que llevamos con nosotros no es un complemento cualquiera: es arte. 
	Shiesty nace con el objetivo de poner en contacto a todos estos coleccionistas con el fin de dar a conocer sus productos a traveés de la venta de estos. Sin embargo, esta no es una venta cualquiera. Shiesty tiene el objetivo de conseguir el mayor valor de venta posible buscando el preic

\capitulo{2}{Objetivos del proyecto}

Este apartado explica de forma precisa y concisa cuales son los objetivos que se persiguen con la realización del proyecto. Se puede distinguir entre los objetivos marcados por los requisitos del software a construir y los objetivos de carácter técnico que plantea a la hora de llevar a la práctica el proyecto.

\capitulo{3}{Conceptos teóricos}

%En aquellos proyectos que necesiten para su comprensión y desarrollo de unos conceptos teóricos de una determinada materia o de un determinado dominio de conocimiento, debe existir un apartado que sintetice dichos conceptos.

%Algunos conceptos teóricos de \LaTeX{} \footnote{Créditos a los proyectos de Álvaro López Cantero: Configurador de Presupuestos y Roberto Izquierdo Amo: PLQuiz}.

%\section{Secciones}

%Las secciones se incluyen con el comando section.

%\subsection{Subsecciones}

%Además de secciones tenemos subsecciones.

%\subsubsection{Subsubsecciones}

%Y subsecciones. 


%\section{Referencias}

%Las referencias se incluyen en el texto usando cite~\cite{wiki:latex}. Para citar webs, artículos o libros~\cite{koza92}, si se desean citar más de uno en el mismo lugar~\cite{bortolot2005, koza92}.


%\section{Imágenes}

%Se pueden incluir imágenes con los comandos standard de \LaTeX, pero esta plantilla dispone de comandos propios como por ejemplo el siguiente:

%\imagen{escudoInfor}{Autómata para una expresión vacía}{.5}



%\section{Listas de items}

%Existen tres posibilidades:

%\begin{itemize}
%	\item primer item.
%	\item segundo item.
%\end{itemize}

%\begin{enumerate}
%	\item primer item.
%	\item segundo item.
%\end{enumerate}

%\begin{description}
%	\item[Primer item] más información sobre el primer item.
%	\item[Segundo item] más información sobre el segundo item.
%\end{description}
	
%begin{itemize}
%\item 
%\end{itemize}

%\section{Tablas}

%Igualmente se pueden usar los comandos específicos de \LaTeX o bien usar alguno de los comandos de la plantilla.

%\tablaSmall{Herramientas y tecnologías utilizadas en cada parte del proyecto}{l c c c c}{herramientasportipodeuso}
%{ \multicolumn{1}{l}{Herramientas} & App AngularJS & API REST & BD & Memoria \\}{ 
%HTML5 & X & & &\\
%CSS3 & X & & &\\
%BOOTSTRAP & X & & &\\
%JavaScript & X & & &\\
%AngularJS & X & & &\\
%Bower & X & & &\\
%PHP & & X & &\\
%Karma + Jasmine & X & & &\\
%Slim framework & & X & &\\
%Idiorm & & X & &\\
%Composer & & X & &\\
%JSON & X & X & &\\
%PhpStorm & X & X & &\\
%MySQL & & & X &\\
%PhpMyAdmin & & & X &\\
%Git + BitBucket & X & X & X & X\\
%Mik\TeX{} & & & & X\\
%\TeX{}Maker & & & & X\\
%Astah & & & & X\\
%Balsamiq Mockups & X & & &\\
%VersionOne & X & X & X & X\\
%}

	En este apartado del informe trataré de abordar los puntos teóricos necesarios para afrontar de manera más clara el resto de secciones. La selección de los conceptos teóricos incluidos en esta sección es puramente personal. El objetivo de este apartado es resumir puntos clave que me han servido a mí para conseguir los objetivos marcados en el punto anterior.
	
\section{Definición de subasta y tipos de subastas}

	El término subasta se define como la situación donde dos o más compradores pujan distintas cantidades de dinero por un mismo producto, siendo uno de ellos el comprador final. Para llevarla a cabo se deben seguir una serie de reglas y directrices marcadas antes del comienzo de la subasta. El objetivo de este tipo de ventas es conseguir una mejor venta posible, llegando a un acuerdo justo entre vendedor y comprador.
	
	Según \cite{apuntes:ecommerce_adwords}, actualmente existen muchos tipos de subasta, pero todas ellas se basan en cuatro pilares fundamentales o tipos:
	
	\begin{description}
		\item[Precio descendente u holandesa:] la subasta comienza con un precio elevado y va disminuyendo a medida que nadie interrumpe al subastador. En el momento que alguien quiere pujar, para al subastador en el precio que desee. Este hombre será el ganador de la subasta.
		\item[Precio ascendente o inglesa:] la subasta comienza con un precio mínimo y va subiendo a medida que los posibles compradores lanzan pujas de igual o mayor precio de lo que el subastador indica. El ganador de la subasta será aquella persona que lance la puja más alta.
		\item[De sobre cerrado a primer precio:] los pujadores introducen el precio en un sobre cerrado. Una vez estén todos listos, se abren los sobres y gana la persona que marcó la puja más alta, teniendo que pagar tal cantidad por lo subastado.
		\item[De sobre cerrado a segundo precio o Vickrey:] la dinámica es igual que el anterior tipo de subasta, pero con una diferencia: la persona con la puja más alta pagará el segundo precio más alto por el producto.
	\end{description}
	
	Una pregunta que puede surgirnos es por qué existen tantos tipos si el objetivo es el mismo. ¿Por diversión? ¿Por cambiar? La respuesta reside en el objetivo de alcanzar el mejor acuerdo posible vendedor-comprador, es decir, sacar el máximo grado de objetividad a la venta. 
	
	En una subasta podemos diferenciar entre las pujas normales y el precio reserva, este último siendo nuestra máxima cantidad a pujar. Referenciando a los tipos explicados hace unas líneas, la subasta de sobre cerrado a primer precio no es una subasta del todo óptima. El pujador no trata de lanzar su precio de reserva y esto puede hacer que, si la subasta se repitiese, el precio se ajuste mucho y el mercado se vuelva volátil. Sin embargo, la subasta de sobre cerrado a segundo precio o Vickrey incita al pujador a adoptar una estrategia dominante y lanzar su puja máxima, obteniendo un grado de venta muy satisfactorio para ambas partes.
	
	En nuestro caso, el tipo de subasta elegido para nuestra aplicación será subasta de precio ascendente o inglesa, sin duda el tipo de subasta más frecuente.
	
	

\section{Definición de Machine Learning, algoritmos y evaluación}

	El Aprendizaje Automático o Machine Learning es una parte de la informática muy ligada al concepto de AI o Inteligencia Artificial. Según \cite{oracle:homepage}, el término IA se refiere a ``aplicaciones que realizan tareas complejas para las que antes eran necesaria la intervención humana''. Principalmente se busca entrenar a la máquina con el uso de algoritmos y datos de forma que esta aprenda de manera gradual. No hay que decir que cuanto mayor sea el número de datos y cuanto mayor sea la calidad de estos, mayor precisión se conseguirá. 
	
	El término Machine Learning podemos atribuírselo a Arthur Samuel quien, en 1952, creó un software capaz de aprender y jugar a las damas. Aún así, este campo ha evolucionado a pasos agigantados, siendo actualmente uno de los puntos más relevantes dentro del campo de la informática. Entre los hitos más significativos, destacan principalmente unos estudiantes de la Universidad de Stanford. Estos alumnos consiguieron desarrollar en 1979 un software capaz de pilotar un carro de manera autónoma sin que este chocara con ningún obstáculo.  No podemos olvidarnos de uno de los momentos más significativos de nuestra época cuando el gran Garri Kaspárov fue derrotado por Deeper Blue en una partida de ajedrez en 1997, siendo la primera vez que el ajedrecista perdía contra una máquina.
	
	Como se ha marcado en esta introducción, esta rama de la informática se centra en desarrollar algoritmos y modelos autosuficientes, es decir, que sean capaces de aprender de manera automática a través de datos y la experiencia. Dentro de esta rama, encontramos infinidad de técnicas y algoritmos, aunque me voy a quedar con los más comunes.
	
	Según \cite{ibm:machine_learning}, el aprendizaje supervisado es una rama del aprendizaje automático donde se entrena a un modelo con una serie de datos que contienen cuál es su entrada y cómo debe ser su salida. Un ejemplo de ello es la venta de un local: el modelo recibe como variables de entrada el número de baños, los metros cuadrados, el año de edificación... y el modelo procesa todo dando como salida el precio del local. Dentro de este tipo de aprendizaje pueden abordarse dos tipos de problemas: regresión si lo que devuelve la máquina es un valor numérico continuo (siguiendo con el ejemplo anterior, el precio del local), o clasificación si lo que devuelve es un valor categórico (clasificar imágenes de vehículos entre coches o camiones).
	
	Son muchos los algoritmos presentes en este grupo, pero los más notables son:
	
	\begin{description}
		\item[Regresión lineal:] es uno de los modelos más simples y fáciles de interpretar. Destaca porque modela la relación entre una variable dependiente (variable que se intenta predecir) y una o varias independientes (variables no resultado que sirven para conseguir la predicción) mediante una recta, consiguiendo así predicciones más precisas y comprender como influyen.
		\item[Regresión logística:] es un algoritmo fundamental a la hora de enfrentarnos a problemas de clasificación binaria. Destaca por utilizar la función logística (función matemática para transformar cualquier valor en otro con rango de 0 a 1) consiguiendo así modelar la probabilidad de que una instancia pertenezca o no a una clase.
		\item[SVM:] se utiliza tanto en regresión como en clasificación. Son muy eficaces en espacios de alta dimensión y casos donde el número de muestras es mucho menor que el número de dimensiones. El objetivo es separar las instancias en el espacio de características a través de la búsqueda del hiperplano, maximizando así el margen entre clases.
		\item[Árboles de decisión:] el objetivo es dividir el espacio de características en regiones y asignar una etiqueta a cada una de ellas. Destacan por ser muy intuitivos y fáciles de interpretar, así como versátiles ya que pueden manejar tanto variables numéricas como categóricas.
		\item[Random forest:] conjunto de árboles de decisión que trabajan de manera combinada para alcanzar una mayor precisión y así evitar el sobreajuste. Destacan porque ayudan a mejorar la precisión de las predicciones a partir de promediar los resultados de un conjunto de árboles de decisión independientes. Entre sus técnicas destaca el \emph{bagging}, consiguiendo así reducir la varianza del modelo y mejorar la generalización.
	\end{description}

%	Por otro lado estaría el aprendizaje no supervisado donde los datos no tienen ni etiquetas ni salidas asociadas. Su objetivo es que la máquina aprenda a base de identificar patrones en los datos y extraer información útil de ellos. Para explicar mejor esto, pongamos un ejemplo: teniendo un conjunto de transacciones bancarias, determinar grupos de personas según las similitudes de compras. Dentro de ello podemos distinguir tres tipos de problemas:
%	
%	\begin{description}
%		\item[Clustering:] el objetivo es agrupar los datos en conjuntos o clusters según existan similitudes entre ellos. Por indicar alguno de los algoritmos dentro de este tipo, destacan:
%			\begin{description}
%				\item [K-Means:] divide los datos en k grupos teniendo en cuenta la minimización de la varianza. El objetivo es que los centros de estos grupos se vayan actualizando de manera iterativa minimizando así el sumatorio de las distancias cuadradas entre los datos y el centro más cercano.
%				\item [Clustering Jerárquico:] consiste en agrupar datos a través de la fusión iterativa de los clusteres más similares. El resultado suele ser representado con un dendrograma.
%				\item [Density-Based Spatial Clustering of Applications with Noise (DBSCAN):] agrupa los datos en función de la densidad. Destaca por ser un algoritmo de gran ayuda a la hora de detectar ruido.
%			\end{description}
%		\item[Reducción de dimensionalidad:] el objetivo es reducir el número de variables en un conjunto de datos sin perder la mayor parte de la información. Ejemplos de algoritmos son:
%			\begin{description}
%				\item [Análisis de Componentes Principales (PCA):] reduce dimensionalidad transformando los datos en nuevas variables ordenadas según al cantidad de varianza que explican.
%				\item [T-Distributed Stochastic Neighbor Embedding (t-SNE):] reduce la dimensionalidad disminuyendo el espacio y manteniendo las relaciones de proximidad locales. Muy útil para la búsqueda de patrones.
%			\end{description}
%		\item[Asociación:] el objetivo de la máquina es que aprenda a base de relacionar variables en un conjunto de datos determinado. Un ejemplo dominante dentro de este grupo son la sección de recomendación de webs de compras. Ejemplos de algoritmos dentro de ello son:
%			\begin{description}
%				\item [Equivalence Class Clustering and bottom-up Lattice Traversal (Eclat):] permite encontrar conjuntos de datos frecuentes en bases de datos transaccionales.
%				\item [Frequent Pattern Growth (FP-Growth):] mismo objetivo que el anterior pero la dinámica consiste en la contrucción de un árbol de patrones frecuentes.
%			\end{description}
%	\end{description}
%
%	Si combinamos las ideas fundamentales de estos dos aprendizajes anteriores, damos lugar al aprendizaje semisupervisado. Simplemente permite que la máquina aprenda a base de datos sin etiqueta y con ella.
%	
%	Otro tipo de Machine Learning es el conocido como aprendizaje por refuerzo. La idea fundamental de este tipo de aprendizaje es que la respuesta del modelo mejore a partir de una retroalimentación en forma de recompensas. En otras palabras, el algoritmo va a aprender basándose en lo que le rodea y siguiendo la filosofía “ensayo-error”. Se utiliza en infinidad de sectores y los algoritmos por excelencia dentro de este aprendizaje son:
%			\begin{description}
%				\item [Q-Learning:] actualiza los valores teniendo en cuenta las recompensas recibidas y las estimaciones futuras.
%				\item [State-Action-Reward-State-Action (SARSA):] actualiza los valores teniendo en cuenta las acciones presentes y las futuras.
%				\item [Algoritmo de Policy Gradient:] actualiza los valores según la dirección del gradiente de la recompensa esperada.
%			\end{description}
%	
%	Por último, otro aprendizaje muy presente en diversos sectores es el Deep Learning. El objetivo de este campo es emular el comportamiento de un cerebro humano a partir de algoritmos de redes neuronales artificiales formadas por múltiples capas de procesamiento capaces de representar datos de alto nivel. Ejemplos muy característicos son el reconocimiento de imágenes, detección de objetos, detección de fraudes, traducción automática...
	
	En nuestro caso, nuestro problema responde a un aprendizaje supervisado. Aunque lo veremos más adelante, se han probado dos algoritmos distintos, siendo Regresión Lineal el definitivo.
	
\section{Arquitectura de aplicaciones móviles}

	Cuando hablamos de la arquitectura de aplicaciones móviles hablamos tanto de los componentes y su interacción como de la estructura y diseño que albergan dentro de ellas. Para explicarlo de una manera óptima, podemos dividir esta estructura en una serie de capas:
	
	\begin{description}
		\item[Capa visual o Frontend:] esta capa recoge todo lo relacionado con la interfaz y la experiencia de usuario (UI y UX respectivamente). Sus componentes principales (siguiendo el patrón MVC explicado en el punto 4 de este informe: Técnicas y herramientas) son:
			\begin{enumerate}
				\item Controladores : albergan la lógica de presentación de la aplicación.
				\item Vistas: muestran al usuario la información que les llega desde los controladores.
			\end{enumerate}
		\item[Capa lógica]: dentro de esta capa se encuentran los modelos según el patrón MVC. No soy partidario de meterlo ni en la capa anterior ni en la posterior porque no responde como tal a ninguna de las dos. En este caso, los modelos son el puente entre ambas partes, pues se encargan del manejo de los datos.
		\item[Capa de datos o Backend:] capa que tiene como objetivo almacenar datos y prepararlos para su posterior uso en la aplicación. Dentro de esta capa encontramos componentes como:
			\begin{enumerate}
				\item APIs : permiten hacer llamadas a servicios externos con el fin de recuperar datos.
				\item Bases de datos: almacenan datos de manera persistente.
			\end{enumerate}
	\end{description}
	
	En nuestro caso, Flutter será quién protagonice la capa visual y la capa lógica, mientras que Firestore Database será quien marque la capa de datos.
	
\section{Bases de datos en la Nube}

	Como hemos apuntado en la sección anterior, una base de datos es un conjunto de datos almacenados y que tienen relación entre sí. Los SGBD o Sistemas Gestores de Bases de datos son los programas que gestionan todo esto. Dentro de las bases de datos, podemos diferenciar dos grandes tipos:
	
	\begin{enumerate}
		\item Bases de datos SQL o relacionales
		\item Bases de datos NoSQL o no relacionales
	\end{enumerate}
	
	Las bases de datos SQL almacenan una serie de datos y sus relaciones con otras tablas a través de claves, datos que son los mismos en tablas distintas. Cada registro se ve representado por filas mientras que las columnas representan los campos de cada registro. Esta filosofía de almacenamiento nos permite acceder de manera sencilla a cualquier dato. Para ello, contamos con el lenguaje SQL.
	
	Por otro lado se encuentran las bases de datos NoSQL, las cuales surgen por la necesidad de almacenar la información sin seguir una estructura determinada. Se utilizan cuando el número de datos a albergar es muy grande y necesitamos rapidez en el tratamiento de estos datos, cosa que las anteriores, en general, no permiten.
	
	Siguiendo con nuestro ejemplo, la base de datos a utilizar es Firestore Database, base de datos NoSQL.
	
\capitulo{4}{Técnicas y herramientas}

%Esta parte de la memoria tiene como objetivo presentar las técnicas metodológicas y las herramientas de desarrollo que se han utilizado para llevar a cabo el proyecto. Si se han estudiado diferentes alternativas de metodologías, herramientas, bibliotecas se puede hacer un resumen de los aspectos más destacados de cada alternativa, incluyendo comparativas entre las distintas opciones y una justificación de las elecciones realizadas. 
%No se pretende que este apartado se convierta en un capítulo de un libro dedicado a cada una de las alternativas, sino comentar los aspectos más destacados de cada opción, con un repaso somero a los fundamentos esenciales y referencias bibliográficas para que el lector pueda ampliar su conocimiento sobre el tema.

	Habiendo definido los aspectos teóricos más clave para entender un poco mejor el trabajo, se define a continuación cuáles han sido las técnicas y herramientas utilizadas para alcanzar nuestro objetivo. Aunque las hemos nombrado en alguno de los apartados anteriores, esta sección del informe explica de manera más particular qué nos aportan y cuáles son las ventajas frente a otras.
	
\section{Flutter}

	Flutter es un SDK (Software Development Kit) de código abierto desarrollado por Google. Su objetivo principal es la creación de aplicaciones multiplataforma con un único código base. Se basa en el lenguaje de programación Dart, el cual es un lenguaje de programación orientado a objetos desarrollado por Google con el fin de crear aplicaciones web y de escritorio de una manera rápida y sencilla.
	
	Volviendo al framework, Flutter es una herramienta que permite crear interfaces de usuario atractivas visualmente y de una alta personalización. Hay que añadir que dispone de una amplia gama de widgets predefinidos y listos para ser utilizados por cualquier usuario.
	
	En cuanto a su rendimiento, Flutter destaca por poseer su propio motor de renderizado eliminando así la necesidad de dependencia de componentes de interfaz de usuarios nativos.
	
	Aún que son muchas las ventajas  que presenta esta herramienta, hay que destacar una frente a las demás: la posibilidad de crear multiplataforma. Flutter nos va a permitir reutilizar el mismo código para crear nuestra aplicación en las plataformas iOS, Android, web, Windows, MacOS y Linux. A esta ventaja, añadir la posibilidad de recargar el código en caliente, aspecto muy demandado por cualquier persona que se dedique a la programación.
	
	Para explicarlo de una manera sencilla, la unidad mínima de este lenguaje es el widget. Estos objetos son los que van a conformar cualquiera de las vistas que tenga nuestra aplicación. Estos widgets pueden estar formados a su vez por dos o más widgets, consiguiendo una libertad de edición y reutilización de widgets enorme. A esto añadir la gran cantidad de librerías con widgets ya predefinidos y listos para ser usados.
	
	En lo que a este trabajo respecta, no todo son ventajas. La curva de aprendizaje de Flutter tiene muy poca pendiente y esto hace que conseguir una aplicación base llegue a ser frustrante. Solo para poder descargar el entorno debes tener unos requisitos e instalar una serie de programas previos que hacen que la probabilidad de error sea elevada. A esto hay que añadir que desde los últimos años, debido a la aparición de otros competidores, Flutter ha dejado de ser protagonista y su comunidad y documentación a disminuido de manera drástica, llegando a salir noticias actuales como el despido de gran parte del equipo de Flutter en Google.
	
	
\section{Firestore database}

	Siguiendo un poco la herramienta utilizada en la capa visual, buscamos qué Sistema de Gestión de Base de Datos podría ser el óptimo para trabajar con Flutter. La verdad que encontrarlo resuelto sencillo pues Firestore Database también está desarrollado por Google y ofrece facilidades siempre que lo conectes a un SDK como Flutter.
	
	Firestore Database es una base de datos alojada en la nube que responde al tipo NoSQL. Entre sus características principales destacan la rapidez de respuesta y su sincronización de datos en tiempo real. La información se organiza en colecciones y documentos, los cuales podríamos comparar con tablas y registros respectivamente. Sin embargo, es una base de datos que no requiere de esquemas predefinidos, permitiendo así cambios en la estructura de los datos.
	
	Si tuviésemos que definir alguna ventaja más, podríamos afirmar que Firestore es capaz de gestionar grana cantidad de datos y sincronizarlos de manera rápida, así como una gran seguridad a la hora de autenticarnos. Sin embargo, existen algún que otra desventaja difícil de evitar durante este trabajo.
	
	Nunca había tratado directamente con una base de datos distinta al formato SQL, lo que provocó que mi forma de definir las tablas y tratar los datos no fuese del todo correcta en varios momentos críticos del trabajo. Por otro lado, la conexión con la base de datos es una etapa con muchos pequeños pasos intermedios que, al igual que la preparación de Flutter, pueden desembocar en errores si no se cumplen al pie de la letra.


\section{Machine Learning con Visual Studio Code}

	Aunque existen otras formas de realizarlo, la parte de preparación de los datos y entrenamiento del modelo lo he realizado con Visual Studio Code. En otras asignaturas de esta carrera hemos utilizado Anaconda para trabajar sobre Jupyter Notebook, pero gracias a una extensión he trabajado de manera más sencilla así. 
	
	Por otro lado, tanto el script para lanzar ejecutar el modelo y la API han sido realizados con la misma herramienta pero con extensión .py	
	
\section{Render}

	Render es un sitio web que permite desplegar infinidad de aplicaciones web, servicios backend, bases de datos... de una manera sencilla. Aunque hablaremos más a fondo en otros apartados de cómo ha sido la conexión, me gustaría destacar dos ventajas frente a otros competidores como Fly o Heroku:
	\begin{enumerate}
		\item Es gratuita. Las otras parecen ofrecer servicios gratis pero acaban pidiendo un método de pago.
		\item Si tenemos conectado nuestro espacio a un repositorio de Git, se realizará un autodeploy cuando esta detecte algún commit en la carpeta indicada.
	\end{enumerate}

\capitulo{5}{Aspectos relevantes del desarrollo del proyecto}

%\begin{enumerate}
%\def\labelenumi{\arabic{enumi}.}
%\tightlist
%\item
%  \textbf{Travis}: realizaba una compilación del proyecto, ejecutaba los
%  test unitarios, ejecutaba \emph{Lint}, ponía en marcha un emulador de
%  Android, y ejecutaba los Android test sobre este. Al finalizar,
%  enviaba los resultados a Codecov y SonarQube.
%\item
%  \textbf{Codecov}: realizaba un análisis sobre la cobertura de los test
%  unitarios.
%\item
%  \textbf{CodeClimate}: ejecutaba cuatro motores de chequeo
%  (\emph{checkstyle}, \emph{fixme}, \emph{pmd} y \emph{markdownlint})
%  sobre el código para detectar posibles problemas o vulnerabilidades en
%  él.
%\item
%  \textbf{SonarQube}: analizaba código duplicado, violaciones de
%  estándares, cobertura de tests unitarios, \emph{bugs} potenciales, etc.
%\item
%  \textbf{VersionEye}: chequeaba todas las dependencias utilizadas en la
%  aplicación y comprobaba si estaban actualizadas, si tenían algún
%  problema de seguridad conocido, o si violaban la licencia del
%  proyecto.
%\end{enumerate}

	Con todo introducido, pasamos a desarrollar uno de los puntos clave de la memoria de este trabajo: explicar cómo y por qué se ha realizado el trabajo de esta forma. Aunque es cierto que se podría explicar todo el proceso seguido, creo que es conveniente centrarnos en los aspectos que han sido claves para llegar al proceso final y que, muchos de ellos, han sido verdaderos quebraderos de cabeza.
	
	Si algo define este trabajo es la capacidad de una persona para conseguir realizar una aplicación y un modelo de Machine Learning comenzando completamente de cero, pues partía de no conocer absolutamente nada de lo que hoy es el resultado final. 
	
	A modo de ser ordenado, vamos a dividir este apartado en tres partes clave:
	\begin{itemize}
		\item Preparación del dataset.
		\item Entrenamiento del modelo y conexión con aplicación.
		\item Estructura general de la aplicación y control de estados.
	\end{itemize}
	
\section{Creación y conexión del modelo}
	
	Como apuntaba al principio de este apartado, mis conocimientos en el ámbito del Machine Learning eran ínfimos. Había leído algún artículo relacionado con ello, pero nunca había entrado en detalle. Por esta razón, la parte de creación del modelo de predicción de precios fue la primera en realizarse. 
	
	Tras hablar con mis tutores como podríamos afrontarla, marcamos como primer objetivo la búsqueda de un dataset que respondiera de manera eficaz a nuestras necesidades. Fueron muchos los que encontramos, teniendo como palabras clave de búsqueda "E-commerce" o "Watches", entre otras. La búsqueda no fue sencilla pues mi desconocimiento de la matería y la gran cantidad de dataset presentes en la red no me permitía elegir con precisión cual podía ser un buen punto de partida.
	
	Tras varios días de búsqueda, el sitio web "Kaggle" me brindó el dataset definitivo: "Luxury Watch Listings". Este dataset contenía nada más y nada menos que más de 280000 relojes con sus características y su precio de venta. Tras ponerlo en común, fue el elegido.
	
	Una vez elegido, busqué toda información necesaria para tratarlo y poder crear un modelo que devolviese una predicción del precio de cualquier reloj. Leyendo en diversas fuentes de información pude darme cuenta que el problema principal a la hora de entrenar un modelo es la presencia de valores nulos. Concretamente, nuestro dataset poseía los datos nulos indicados en la Tabla 5.1.
	
\tablaSmall{Campos vacíos en los datos}{|c|c|c|c|}{campos-vacios}{
  Nombre & Campos vacíos & Nombre & Campos vacíos \\
  \midrule
  Unnamed & 0 & name & 72585 \\
  price & 406 & brand & 131 \\
  model & 30466 & ref & 43152 \\
  mvmt & 196685 & casem & 164271 \\
  bracem & 174896 & yop & 134 \\
  cond & 75987 & sex & 95805 \\
  size & 23597 & condition & 212922 \\
}

	Aparentemente no eran muchos los campos nulos en proporción a las 280000 lineas del dataset. Sin embargo, esto no fue del todo cierto. Revisando previamente el archivo, vi que existian múltiples datos que nos iban a dificultar nuestro trabajo. Por esta razón, se procedió a hacer un estudio previo desde la aplicación Microsoft Excel. El proceso consistió en crear una tabla dinámica con todos los datos y encabezados, de forma que se pudiera ver los distintos nombres o números que formaban el rango de datos. 
	
	Tras estudiar los datos y saber que la siguiente fase del proceso era convertir todo esa información en números para poder entrenar el modelo, se llegó a una serie de soluciones a problemas que se plantean a continuación:
	
\begin{description}
	\item[Columna unnamed:] esta columna, a primera vista, parecía que marcaba de manera incremental el número de filas del archivo. Sin embargo, cuando cambiabamos de marca de reloj, volvía a comenzar en uno. Por tanto, trataba de marcar el número de relojes por marca. Resumiendo, fue inútil trabajar con ella para conseguir nuestro objetivo. La solución que se planteó fue el borrado de esta.
	\item[Columna name:] sin duda fue la columna más problemática de todo el archivo. Más que ser el nombre del reloj, era una descripción de este. El problema principal fue que la descripción no seguía ninguna estructura: marcaba los datos que quería, como quería, repetía datos de otras columnas... En resumen, ha sido y es una columna intratable. Por tanto, la solución fue su borrado.
	\item[Columna price:] esta columna no presentaba apenas errores. Simplemente marcaba el precio de cada reloj. Sin embargo, aquel producto que no se le fuese marcado un precio, adquería el valor “Price on request” o directamente se encontraba en blanco. La solución que se planteó fue dar a estas dos excepciones el valor -1 para entender que no existe información relativa sobre ese reloj. Hay que añadir que se eliminó el símbolo dolar para evitar errores.
	\item[Columna brand:] no presentó ningún error. Simplemente, se categorizó para su posterior uso en el modelo a crear. Las celdas sin información pasaron a valer -1.
	\item[Columna model:] no presentó ningún error. Simplemente, se categorizó para su posterior uso en el modelo a crear. Las celdas sin información pasaron a valer -1.
	\item[Columna ref:] siendo como soy, un apasionado de los relojes, sabía que tener la referencia del reloj que buscamos es uno de los puntos más claves a la hora de reconocer el reloj. Sin embargo, la columna de nuestro dataset no se encontraba del todo limpia. Aun así, no quise perder tan valiosa información y me di cuenta de que toda referencia que lleva una letra siempre aparecería en mayúscula. Por tanto, la solución que se dió fue eliminar toda aquella palabra que tuviera al menos una minúscula.
	\item[Columna mvmt:] no presentó ningún error. Simplemente, se categorizó para su posterior uso en el modelo a crear. Las celdas sin información pasaron a valer -1.
	\item[Columna casem:] no presentó ningún error. Simplemente, se categorizó para su posterior uso en el modelo a crear. Las celdas sin información pasaron a valer -1.
	\item[Columna bracem:] no presentó ningún error. Simplemente, se categorizó para su posterior uso en el modelo a crear. Las celdas sin información pasaron a valer -1.
	\item[Columna yop:] la información se brindaba de cuatro formas diferentes:
		\begin{description}
			\item - Un único año (Ej. 2024)
			\item - Varios años separado por comas (Ej. 2021, 2022, 2023)
			\item - Un año seguido de la palabra aproximación (Ej. 2020 (Approximation))
			\item - No se sabe el año y se marca como Unknown, o se deja en blanco.
		\end{description}
	La solución que se aportó fue únicamente marcar un año por celda. Por tanto, para el primer caso no se hizo nada; el segundo caso se hizo la media entre todos los años que se marcaban; para el tercero se quitó la palabra “Approximation”; y para el último se cambió el valor “Unknown” o celda blanca por -1.
	\item[Columna cond:] esta columna tenía los datos perfectamente dispuestos para ser categorizados. Sin embargo, se vió que terminaba a mitad de camino, siendo la columna “condition” quien seguía el trabajo de esta. La solución que se aportó fue combinar ambas columnas en una misma y categorizar, marcando como -1 aquellas celdas que se encontraban vacías.
	\item[Columna sex:] no presentó ningún error. Simplemente, se categorizó para su posterior uso en el modelo a crear. Las celdas sin información pasaron a valer -1.
	\item[Columna size:] debido a que fue un campo que oposu muchas dificultades, la solución que se aportó fue simple: borrar la columna. No es un campo que aportara información útil comparado con otros campos como la marca, el estado, el precio... que son variables que van a definir casi por completo la predicción objetivo.
	\item[Columna condition:] explicada en apartado "Columna cond".
\end{description}
	
	Para llegar a este momento pasaron par de semanas hasta que creí encontrar la solución más óptima para poder empezar a crear y entrenar un modelo. Realicé dos líneas de entrenamiento: Regresion Lineal y Random-Forest-Regresor. Los resultados de las métricas fueron los indicados en la Tabla 5.2.
	
\tablaSmall{Comparación de métricas}{|c|c|c|}{comparacion-metricas}{
   & Regresion Lineal & Random-Forest-Regresor \\
  \midrule
  Mean-Squared-Error & 2033534792.8564095 & 1380488875.0762675 \\
  Root-Mean-Squared-Error & 45094.73132037056 & 37154.930696695796 \\
  R2-score & 0.4078856446303635 & 0.598036245442073 \\
  R2-score-ajustado & 0.4078464450275924 & 0.5980096343333164 \\
  Mean-absolute-percentage-error & 1.2132659733005473 & 0.23691091500423903 \\
}
	
	Tras una reunión vimos que la métrica "Mean Absolute Percentage error" era muy buena, mientras que las demás no lo eran tanto. Esto nos hizo dudar de cómo se estaban tratando, llegando tras varios días a la conclusión de que había cuatro errores clave en la forma de tratar mis datos:
	
	\begin{description}
		\item - Al categorizar por mi cuenta, estaba queriendo decir al modelo de alguna forma que la marca de relojes que recibia el valor 1 era peor que la marca de relojes que recibía el valor 2, y esto no tenía por qué ser así.
		\item - Dar el valor -1 a aquellos campos nulos o sin información relevante no fue una buena práctica para entrenar al modelo.
		\item - Durante la realización de esta parte, la consola me devolvía error con el tipo de datos. Tras dos semanas intentado ver que ocurría vi que estaba cometiendo un error: los valores int no tienen soporte directo para representar NaN, sin embargo float sí. Esto me hizo tener que volver a empezar de nuevo y pensar de manera distinta el proceso.
		\item - La libreria pandas trata los valores faltantes de diferente forma que la libreria numpy, siendo pandas pd.NA y numpy np.NaN.
	\end{description}
	
	En conclusión, tuvimos que cambiar de nuevo la forma de trabajar. Dicen que a veces es más fácil desaprender y volver a realizar las cosas de nuevo. En mi caso, me surgió efecto.
	
		Lo primero que hice fue ver como podía limpiar los datos, cosa que reutilicé parte de lo que ya había pensado anteriormente. No traté de categorizar, simplemente me quedé con lo que me interesaba y lo que no lo vacié. Con esto listo, pensé en como podía completar los campos vacíos: 
		\begin{description}
		\item - Para la columna precio, intenté aplicar la media de aquellos que comparten marca y modelo. No sirvió, pues existen relojes de la misma marca y modelo que se diferencian mucho, bien por la correa, la edición, el año...
		\item - No podía dar valores porque sí al sexo, brazalete, caja, año...
	\end{description}
	Debido a no tener forma de dar valores a los campos faltantes, me salieron bastantes campos sin información. Sin embargo, empecé a probar combinaciones de columnas para ver con cuantas me podía quedar sin perder parte de las filas del dataset. Tras varias pruebas, las candidatas fueron:
		\begin{description}
		\item - brand
		\item - model
		\item - condition
		\item - yop
		\item - price
	\end{description}
	
	El resultado final fue un dataset de 160000 líneas con tres columnas que pasarían a ser categóricas en un futuro, una columna númerica y una columna objetivo.
	
		
\section{Entrenamiento del modelo y conexión con aplicación}

	Referenciando al apartado anterior, nos situamos en el dataset final. Tras semanas de trabajo y reuniones con tutores, vimos que debíamos centrarnos en cómo ibamos a comunicar el modelo de Machine Learning con la aplicación de Flutter. Este proceso fue uno de los más duros, pues la información que había en la red era escasa y diferente. Si tuviera que definir este apartado en una frase sería: prueba y error.
	
	Con el dataset preparado, descubrí OneHotEncoder (explicado en parte teórica) y entrené al modelo, consiguiendo un modelo que predecía el precio del reloj. Tras revisar las pocas fuentes de información disponibles, vie que existía una herramienta para conectar el modelo: tflite. Al parecer, todo era muy sencillo. Simplemente debía conseguir guardar el modelo con la extensión .h5, agregar una serie de dependencias en mi programa y descargar algún que otro paquete. Sin embargo, esto o fue así. Cada vez que lanzaba de alguna manera peticiones al modelo, el problema con los tipos de datos afloraba.
	
	Esto no fue el único problema. Un error mío fue entrenar siempre por completo el modelo cada vez que intentaba una cosa nueva. Esto me suponía cerca de tres cuartos de hora esperando a que terminara, muchas veces con el resultado final de no poder ni guardar el modelo. 
	
	Al igual que con el dataset, dar con el problema fue un mundo. Traté de cambiar el tipo de datos, traté hacer conversiones extensiones sobre el modelo... pero nada. El modelo no respondía. Volviendo un poco al primer apartado, pensé que podía ser problema de las versiones de "tensorflow", "Flask"... y así fue: cada vez que realizaba un "pip install tensorflow" se decargaba por detrás una librería "tensorflow-intel" que descuadraba la compatibilidad de unas librerías con otras. Resumiendo, había que buscar otra forma.
	
	Ojeando par de vídeos y leyendo en algún sitio web, ví que podía conseguir mi objetivo dando algo más de vuelta. Como veía que entrenar mi modelo suponía una gran cantidad de tiempo, tomé el consejo de mis tutores y traté de reducir el número de líneas de entrenamiento ya que mi objetivo era que se conectase a la aplicación como fuera.
	
	La extensión elegida fue ".joblib". Una vez conseguida, realicé un pequeño script en python para poder lanzar en local. Con todo ello, conseguí que el script funcionase, pero tuve problemas con la forma de lanzar peticiones desde la terminal. Esto fue un poco frustante, ya que veía como siguiendo la información disponible en red funcionaba y yo no lo conseguía.
	
	No podía destinar el tiempo del que no disponía en ello, así que intenté trabajar con Postman, herramienta muy intuitiva que yo había utilizado tanto en el mundo estudiantil como en el mundo laboral. Tras un tiempo configurando todo y corrigiendo errores, conseguí dar respuesta a mi petición POST, cosa que celebré como un gol en el ambito futbolístico.
	
	La siguiente fase era desplegarlo de alguna forma para que mi aplicación Flutter pudiera hacer la llamada que había conseguido con Postman. Muchos vídeos me recomendaban Fly.io y Heroku. Ambas presumían de ser entornos gratuitos, pero a la hora de la verdad, te obligaban a ingresar un método de pago. Aun así, comenté mi problema con algún compañero y me recomendó Render, el cual sería mi opción final.
	
	El camino de Render era aparentemente sencillo: conectar con mi repositorio de Git y desplegar. Como marqué al principio del apartado, el proceso fue prueba y error. Tras varios commits, horas de intentos y pruebas, dí con los cuatro archivos necesarios para conseguir desplegar la API que haría las peticiones al modelo: Procfile.txt, requirements.txt, el modelo con extensión .joblib y el script con extensión .py.
	
	Con todo listo, la predicción y conexión con mi aplicación estaban listas.
	
\section{Estructura general de la aplicación y control de estados}
	
	Como se ha indicado en el punto 4 de este informe, se ha intentado aplicar el patrón MVC en el trabajo. Al tener tres tablas en nuestra base de datos, es necesario crear tres modelos en nuestro programa. Estos han recibido los nombres de:
	\begin{enumerate}
		\item user.dart (usuario)
		\item watch.dart (reloj)
		\item auction.dart (subasta)
	\end{enumerate}

	Los tres archivos siguen la misma filosofía:
	\begin{enumerate}
		\item Se declaran cada uno de los objetos.
		\item Se estipula cómo se recibe y envía información a la base de datos.
		\item Se declara cualquier función que trate los datos.
	\end{enumerate}
	
	Por otro lado, se han definido las vistas y los controladores en el mismo script ya que Flutter requiere de ello. Como indicábamos en otros apartados, Flutter trabaja con widgets y estos necesitan que las funciones sean definidas para recargarlas cada vez que se producen cambios por diversas acciones, aunque esto lo hablaremos con más detalle en el apartado del manual de usuario.
	
	Un aspecto importante de este trabajo es controlar muy bien los estados de cada reloj y subastas tras realizar una acción cualquiera. Por ello, se enumera a continuación todos los casos tenidos en cuenta:
	\begin{enumerate}
		\item Cuando un usuario crea un reloj, no puede crear dos relojes con un mismo watch nickname ya que luego van a depender de este campo a la hora de crear una subasta.
		\item Una vez creado el reloj, el estado del reloj pasará a estar subido, lo que indicará que ni se encuentra en subasta ni nunca ha sido partícipe de una. Por tanto, el usuario podrá editar y/o eliminar el reloj sin ningún problema.
		\item Cuando un usuario quiera crear una subasta, deberá crearla con un nombre de los relojes asociados a su cuenta. En el momento de su creación, el estado de la subasta pasará a activa y el estado del reloj pasará a estar en subasta. Por tanto, se deben bloquear los botones de editar y eliminar del reloj. En cuanto a la subasta, el creador de la subasta tendrá bloqueados los botones tanto de compra directa como el botón de puja, mientras que el botón de eliminar estará disponible siempre que no haya aplicado nadie a la subasta. En cuanto a otros usuarios distintos al creador, tendrán disponibles los botones de manera contraria a este.
		\item Si el creador de una subasta elimina esta, el estado de la subasta pasará a finalizado y el estado del reloj pasará a disponible, volviendo a habilitar de nuevo los botones de edición y eliminación del reloj.
		\item Si por la razón que sea la subasta terminase, el estado de esta pasaría también ha finalizado, pero el estado del reloj pasaría a vendido y los botones seguirían deshabilitados.
	\end{enumerate}
	
	Estos no han sido los únicos condicionantes a la hora de realizar acciones. Uno muy importante ha sido controlar bien las pujas de los usuarios. Siempre que usuario iba a realizar una puja o la compra directa del reloj, se han tenido en cuenta dos condicionantes:
	\begin{enumerate}
		\item El monedero del usuario debe ser mayor que la cantidad a pujar o la cantidad de compra directa.
		\item El monedero del usuario debe ser mayor al sumatorio de todas las pujas realizadas por un usuario.
	\end{enumerate}
	
	Refiriéndome al último condicionante, nos ponemos en el supuesto de que un usuario gane todas las pujas que aplica. Si ganase todas, debe poder pagar todas. De ahí que se haga la comprobación y, si no pudiese hacer frente a ello, se le avise con una alerta informativa explicando el porqué de todo.

\capitulo{6}{Trabajos relacionados}

	Como tal, no he encontrado un trabajo que sea estrictamente igual a este. Sin embargo, tras hablar y recibir ideas de mis tutores, sí que hemos encontrado algunos que podrían compartir características con Kairos.
	
	Uno de los sitios web por excelencia es Ebay, el cual es uno de los mercados en línea más grande. Ebay consiste en un comercio electrónico global donde los usuarios suben sus productos y pueden vendérselos a otros bien a través de subasta o por venta directa. En su caso, no se centra en un único producto, si no que deja al usuario la libertad de vender lo que él decida. Otra aplicación muy parecida es Wallapop, quién sigue la filosofía de Ebay pero no incorpora la posibilidad de vender los productos a través de subasta.
	
	La diferencia de las dos aplicaciones comentadas antes es que no se centran en un producto como los relojes, lo que hace que tengan de todo pero no se especialicen en algo en concreto. Por eso nace Kairos, para que las personas ligadas al mundo de los relojes encuentren a Kairos un lugar cómodo para poder comprar y vender lo que más les gusta.
	
	Por otra parte, este trabajo se ha realizado utilizando Flutter y Firestore Database. En cuanto a Flutter, existen muchos Trabajos de Fin de Grado que han utilizado esta herramienta, donde cabe destacar el siguiente relizado por D. Daniel Martín Cubero, a quien referencio y felicito por su trabajo: \cite{daniel:dasell}.
	
	Por último, destacar el trabajo de D. Cesar Linares Traseira, quien creó un modelo de predicción de precios de acciones de bolsa a través de redes nuronales siguiendo un modelo de aprendizaje supervisado. A continuación, referencio a este y felicito su trabajo: \cite{cesar:modelo}.
	
	 

\capitulo{7}{Conclusiones y Líneas de trabajo futuras}

%Todo proyecto debe incluir las conclusiones que se derivan de su desarrollo. Éstas pueden ser de diferente índole, dependiendo de la tipología del proyecto, pero normalmente van a estar presentes un conjunto de conclusiones relacionadas con los resultados del proyecto y un conjunto de conclusiones técnicas. 
%Además, resulta muy útil realizar un informe crítico indicando cómo se puede mejorar el proyecto, o cómo se puede continuar trabajando en la línea del proyecto realizado. 

	Con todos los apartados anteriores explicados, llega el momento de realizar una valoración global del trabajo realizado. Es difícil valorarse de manera objetiva a uno mismo, pero creo que una actividad muy efectiva para seguir progresando tanto laboralmente como psicológicamente.
	
	Como marcaba al principio de este informe, el trabajo realizado ha sido un verdadero quebradero de cabeza. En muchas ocasiones lo he comparado con la acción de intentar tirar una pared a cabezazos. Con esto último me quiero referir principalmente a lo difícil que ha sido para mí afrontar una aplicación de este nivel.
	
	Es cierto que las herramientas utilizadas permiten realizar infinidad de cosas increíbles, pero creo que la curva de aprendizaje que carga a su espalda apenas tiene pendiente. Aun así, a base de esfuerzo y dedicación, se ha llegado a crear una aplicación operativa que responde a la mayoría de los objetivos definidos en este informe.
	
	El proyecto puede ser escalado de mil formas, pues existen una gran cantidad de nuevas funciones que podría añadir esta aplicación. Si tuviera que marcar alguna en concreto, me quedaría con:
	
	\begin{enumerate}
		\item Inserción de nuevos tipos de subasta: aunque hemos definido la más común y creo firmemente que sería la más utilizada si hubiese más, podríamos tener en cuenta los tres tiposd e subastas restantes explicados en la parte teórica de este informe.
		\item Conexión con el entorno Google: muchas aplicaciones y sitios web permiten registrarse con la cuenta de Google, cosa que podría incluirse para aumentar el nivel de profesionalidad de nuestra aplicación.
		\item Diseños más profesionales: la aplicación mostrada es una base de una aplicación que podría incluir diseños y efectos de cambios de vista más espectaculares, dando una experiencia de usuario mucho más satisfactoria.
		\item Perfeccionamiento del modelo: este trabajo puede llevarse la mayor cantidad de tiempo a invertir, pues es fundamental que la predicción de precios sea lo más perfecta posible.
	\end{enumerate}
	
	En conclusión, las líneas de mejora de esta aplicación son infinitas y tengo claro que se irán alcanzando con el aprendizaje tanto en mi mundo laboral como en el estudiantil, mundos que son un punto y seguido tras la finalización de este trabajo.


\bibliographystyle{plain}
\bibliography{bibliografia}

\end{document}
